% ============================================================================
% chapters/ch1_introduction.tex
% ============================================================================
\chapter{Introduction}
\label{ch:introduction}

% ----------------------------------------------------------------------------
\section{Problem Statement}
\label{sec:problem}
% ----------------------------------------------------------------------------

Human Activity Recognition (HAR) using wearable inertial measurement units
(IMUs) has progressed from a laboratory research problem to a practical tool
for continuous health monitoring.
In the context of anxiety detection, wrist-worn accelerometer and gyroscope
data can be used to classify physical activities---such as walking, sitting,
or fidgeting---whose patterns correlate with mental health states.
Deep learning architectures, including one-dimensional convolutional neural
networks combined with bidirectional long short-term memory layers
(\cnnbilstm{}), have demonstrated strong classification performance when
evaluated on curated benchmark datasets \citep{chen2021deep}.

However, a significant gap persists between research-grade model evaluation
and reliable operation in production.
The central problem addressed by this thesis is:

\begin{quote}
\textit{Once a HAR model trained on labelled laboratory data is deployed to
process continuous, unlabelled wearable recordings, how can the system detect
that the model's reliability has degraded and adapt without requiring
manual re-labelling?}
\end{quote}

This question is rarely addressed in the HAR literature.
The majority of published work evaluates models under closed-world assumptions:
the training distribution matches the test distribution, all test data are
labelled, and the model is evaluated once rather than continuously.
In a real deployment, none of these assumptions hold.
Users differ in anthropometry, handedness, and movement idiosyncrasies.
Sensor placement may vary between the dominant and non-dominant wrist.
Environmental conditions, firmware updates, and battery degradation introduce
gradual sensor drift.
Over days and weeks, the data distribution shifts away from the training
distribution, and model outputs become progressively less trustworthy---yet no
ground-truth labels exist to detect or quantify this degradation.

The consequences of undetected degradation are particularly severe in a mental
health monitoring context.
Misclassified activity sequences can lead to incorrect anxiety state estimates,
eroding both clinical utility and user trust.
A production HAR system therefore requires not only a performant model, but
also a surrounding infrastructure that
(a)~continuously monitors prediction quality using proxy signals,
(b)~determines when retraining is necessary,
(c)~adapts the model without labelled data, and
(d)~ensures that the entire process is reproducible, versioned, and auditable.

This is fundamentally a Machine Learning Operations (MLOps) problem.
MLOps extends the principles of DevOps---automation, continuous integration,
monitoring, and infrastructure-as-code---to the machine learning lifecycle
\citep{sculley2015hidden}.
While MLOps practices are well-established in domains such as natural language
processing and recommender systems, their application to wearable sensor HAR
remains nascent.

% ----------------------------------------------------------------------------
\section{Research Questions}
\label{sec:rqs}
% ----------------------------------------------------------------------------

The following research questions guide this thesis:

\begin{description}[leftmargin=2cm, labelwidth=1.5cm]
  \item[RQ\,1] How can model degradation in a deployed HAR system be detected
               without access to ground-truth activity labels?
  \item[RQ\,2] Which proxy metrics correlate most reliably with actual
               classification accuracy in a wrist-worn HAR context?
  \item[RQ\,3] When should an automated system trigger model adaptation or
               full retraining, and how can false triggers be suppressed?
  \item[RQ\,4] How effective is curriculum pseudo-labelling as a self-supervised
               adaptation strategy for HAR model updates?
  \item[RQ\,5] What engineering practices (versioning, containerisation,
               CI/CD) are necessary to make a HAR MLOps pipeline
               reproducible and auditable?
\end{description}

% ----------------------------------------------------------------------------
\section{Contributions}
\label{sec:contributions}
% ----------------------------------------------------------------------------

This thesis makes the following contributions:

\begin{enumerate}
  \item A \textbf{14-stage production pipeline} that decomposes the HAR MLOps
        lifecycle into independently executable, composable stages with typed
        configuration and artefact dataclasses
        (\Cref{ch:methodology,ch:implementation}).

  \item A \textbf{three-layer monitoring framework} (confidence, temporal
        consistency, statistical drift) that operates entirely on unlabelled
        predictions and feeds into a voting-based trigger policy
        (\Cref{sec:monitoring}).

  \item A \textbf{tiered trigger policy engine} with multi-metric vote
        confirmation, cooldown periods, and tiered alerting that suppresses
        false retraining triggers (\Cref{sec:trigger}).

  \item A \textbf{curriculum pseudo-labelling procedure} with Elastic Weight
        Consolidation (EWC) regularisation that adapts a pre-trained model
        to new user data without ground-truth labels
        (\Cref{sec:pseudo_labelling}).

  \item An \textbf{AdaBN+TENT test-time adaptation sequence} with a BN
        statistics snapshot mechanism that prevents TENT from corrupting
        AdaBN's output, and dual rollback gates for safety
        (\Cref{sec:tent_adabn}).

  \item A \textbf{calibrated uncertainty pipeline} combining temperature
        scaling and Monte Carlo (MC) Dropout for both monitoring and active
        learning sample selection (\Cref{sec:calibration}).

  \item Integration of all components into a \textbf{containerised, CI/CD-enabled
        deployment} with Docker, GitHub Actions, MLflow experiment tracking,
        and DVC data versioning (\Cref{ch:implementation}).

  \item An \textbf{empirical engineering log} documenting 21 real problems
        encountered during development --- including a sensor unit mismatch
        that reduced cross-user accuracy to 14.5\,\% and the sequence of fixes
        that recovered it --- as evidence-backed case studies for the
        implementation challenges chapter (\Cref{ch:discussion}).
\end{enumerate}

% ----------------------------------------------------------------------------
\section{Scope and Delimitations}
\label{sec:scope}
% ----------------------------------------------------------------------------

The scope of this thesis is bounded as follows.
The sensor modality is restricted to wrist-worn triaxial accelerometers and
gyroscopes (6-axis IMU) as collected by Garmin smartwatches.
The classification target is an 11-class activity taxonomy derived from an
anxiety-related behavioural protocol comprising the activities listed in
\Cref{tab:activity_classes}.

\begin{table}[ht]
  \centering
  \caption{The eleven activity classes used throughout this thesis.}
  \label{tab:activity_classes}
  \begin{tabular}{clcl}
    \toprule
    \textbf{ID} & \textbf{Activity} & \textbf{ID} & \textbf{Activity} \\
    \midrule
    0  & Ear rubbing        & 6  & Nail biting        \\
    1  & Forehead rubbing   & 7  & Nape rubbing       \\
    2  & Hair pulling       & 8  & Sitting (neutral)  \\
    3  & Hand scratching    & 9  & Smoking            \\
    4  & Hand tapping       & 10 & Standing (neutral) \\
    5  & Knuckles cracking  &    &                    \\
    \bottomrule
  \end{tabular}
\end{table}

The deep learning architecture is a fixed \cnnbilstm{}; architecture search
is explicitly out of scope.
The pipeline is designed for single-user adaptation scenarios; federated or
multi-user aggregation is identified as future work.

The thesis does not collect new labelled data.
The primary dataset comprises 26 recording sessions from the Garmin Decoded
export format.
Labelled training data originate from a prior data collection phase described
in \Cref{sec:dataset}.

% ----------------------------------------------------------------------------
\section{Thesis Structure}
\label{sec:structure}
% ----------------------------------------------------------------------------

The remainder of this thesis is organised as follows.

\Cref{ch:background} reviews the literature across four domains: HAR with deep
learning, MLOps practices for continuous ML systems, domain adaptation for
sensor data, and uncertainty quantification in neural classifiers.

\Cref{ch:methodology} presents the methodology, describing the 14-stage
pipeline architecture, the monitoring and trigger framework, and the adaptation
strategies in detail.

\Cref{ch:implementation} details the implementation, covering the repository
structure, CLI design, CI/CD workflow, experiment tracking, and containerised
deployment.

\Cref{ch:evaluation} reports the experimental results: classification
performance, calibration quality, drift detection accuracy, pseudo-labelling
convergence, and the ablation comparison across adaptation methods.

\Cref{ch:discussion} discusses the findings, reflects on the engineering
challenges encountered, identifies limitations, and outlines directions for
future work.
